% concatenated from Preprint_SDEs_Wireless_Malliavin.tex
\documentclass[11pt]{article}
\usepackage{amssymb}
\usepackage{amsthm}
\usepackage{comment}
\usepackage{amsmath}
\usepackage{amsfonts}
\usepackage{amsbsy}
\usepackage{amstext}
\usepackage{mathrsfs}
\usepackage{dsfont}
\usepackage{graphpap}
\usepackage{graphics,graphicx}
\usepackage[T1]{fontenc}
\usepackage{subcaption}
\usepackage{multicol}
\usepackage{booktabs}
\usepackage{array}
\usepackage{bm}
\usepackage{siunitx}
\usepackage{hyperref}
\usepackage[section]{placeins}
\usepackage{xargs}
\usepackage{algorithm}
\usepackage{algpseudocode}
\usepackage{dirtytalk}
\usepackage[pdftex,dvipsnames]{xcolor}
\usepackage[colorinlistoftodos,
   prependcaption,
   textsize=tiny]{todonotes}
\usepackage{mathtools, cuted}
\usepackage[nameinlink,noabbrev]{cleveref}
\usepackage{lineno}
\usepackage{geometry}
\geometry{margin=1in}

\usepackage{natbib}

% Definición de comandos personalizados
\newcommand{\cqd}{\hfill$\Box$}
\newcommand{\f}{{\mathcal F}}
\newcommand{\IR}{{\mathbb R}}
\newcommand{\R}{{\mathbb R}}
\newcommand{\IN}{{\mathbb N}}
\newcommand{\ind}{\mbox{\Large$\chi$}}
\newcommand{\tor}{{\mathbb T}}
\newcommand{\G}{{\mathbb G}}
\newcommand{\beq}{\begin{equation}}
\newcommand{\eeq}{\end{equation}}
\newcommand{\bal}{\begin{align}}
\newcommand{\eal}{\end{align}}
\newcommand{\beqn}{\begin{equation*}}
\newcommand{\eeqn}{\end{equation*}}
\newcommand{\baln}{\begin{align*}}
\newcommand{\ealn}{\end{align*}}
\newcommand{\tbar}{\bar t}
\newcommand{\xbar}{\bar x}
\newcommand{\ep}{\epsilon}
\newcommand{\Pb}{\mathbb P}
\newcommand{\Rl}{\mathbb R}
\newcommand{\E}{\mathbb{E}}
\newcommand{\tf}{\mathcal{F}}
\newcommand{\hac}{\mathcal{H}}
\newcommand{\hact}{\mathcal{H}_T}
\newcommand{\I}{{\mathbb I}}
\newcommand{\grad}{\nabla}
\newcommand{\sign}{\text{sign}}
\newcommand{\Laplace}{\boldsymbol{\Delta}}

% Definición de teoremas y ambientes
\newtheorem{remark}{Remark}
\newtheorem{definition}{Definition}
\newtheorem{theorem}{Theorem}
\newtheorem{hypothesis}{Hypothesis}
\newtheorem{lemma}{Lemma}
\newtheorem{proposition}{Proposition}
\newtheorem{corollary}{Corollary}

% Colores personalizados
\newcommand{\fer}[1]{{\textcolor{red}{#1}}}
\newcommand{\fran}[1]{{\textcolor{blue}{#1}}}

\DeclareGraphicsExtensions{.pdf,.png,.jpg,.tiff}

% Información del documento
\title{Stochastic Analysis of Fade Duration Using Wiener Chaos Expansion and Malliavin Calculus: Optimal Importance Sampling via Adaptive SGD}

\author{Francisco Delgado-Vences\\
Departament de Matemátiques, Universitat Autónoma de Barcelona\\
Bellaterra (Cerdanyola del Valles), España}

\date{\today}

\begin{document}

\maketitle

\begin{abstract}
Characterizing fade duration in wireless channels is fundamental for designing robust communication systems. Classical approaches—Rice's level-crossing theory and Monte Carlo simulation—lack precision for tail events and are computationally prohibitive for rare-event probability estimation. This paper introduces a rigorous framework combining \emph{Wiener Chaos Expansion} (WCE), \emph{Malliavin Calculus}, and \emph{importance sampling with adaptive weights} to analyze fade duration $Z(T)$ distributions. 

Main contributions include: (i) high-accuracy moment estimation and CCDF characterization via WCE minimizing Monte Carlo variance; (ii) \emph{Markovian projection} reducing infinite-dimensional dynamics to tractable systems ($\dim \leq 3$) for Rayleigh, Rician, and Nakagami models under stated assumptions; (iii) \emph{asymptotically optimal importance sampling weights} derived from Malliavin sensitivities, achieving 839--2516$\times$ variance reductions; (iv) a theoretically grounded and \emph{provably efficient adaptive SGD algorithm} with Robbins-Monro step size schedule for parameter estimation. Numerical experiments validate our approach with relative errors $< 0.5\%$, enabling gradient-based optimization of fade duration statistics even for regimes where $P \sim 10^{-15}$, without requiring $\mathcal{O}(1/P)$ samples, by evaluating sensitivities through analytical Malliavin weights.

\end{abstract}

\noindent
\textbf{Keywords:} Fade duration, Malliavin calculus, Wiener chaos expansion, importance sampling, adaptive algorithms, fading channels

\section{Introduction}
\label{sec:introduction}

A fundamental challenge in computational statistics and applied probability is the estimation of occupation times for stochastic processes---that is, the cumulative time a dynamic system spends within a prescribed critical region. For a given stochastic process $X(s)$, the occupation time below a threshold $\gamma$ is formally defined as:
\begin{equation}\label{eq:occupation_time}
    Z(T) = \int_0^T \mathbf{1}_{\{X(s) < \gamma\}} ds,
\end{equation}
where $\mathbf{1}_{\{\cdot\}}$ denotes the indicator function. This discontinuous functional is ubiquitous across applied mathematics and engineering. It appears in financial mathematics for barrier option pricing, in reliability analysis for system failure modeling, and notably in telecommunications for characterizing the fade duration of wireless channels. Understanding the distribution of $Z(T)$---including its rare-event tail probabilities and parameter sensitivities---is critical for robust system optimization.

However, computing the parametric sensitivities (the gradient) of such discontinuous functionals presents a severe numerical bottleneck. Classical gradient estimation techniques, such as Infinitesimal Perturbation Analysis (IPA) or pathwise finite differences, attempt to interchange the derivative and the mathematical expectation. Applying this classical chain rule to Equation \eqref{eq:occupation_time} involves differentiating the indicator function, which yields a Dirac delta singularity. Consequently, evaluating this derivative via standard Monte Carlo simulations is numerically highly unstable, resulting in an estimator with infinite variance. This catastrophic variance inflation systematically destroys the convergence of any gradient-based optimization algorithm, such as Stochastic Gradient Descent (SGD).

To overcome this structural limitation, we draw a methodological parallel with advancements in quantitative finance, where similar discontinuities arise in the payoffs of digital options \cite{fournie1999applications}. The mathematical resolution lies within the framework of Malliavin calculus. By utilizing the Malliavin Integration by Parts Formula (IPF) on the Wiener space, we can analytically bypass the discontinuity of the indicator function. The IPF allows us to transfer the differential operator away from the singular payoff and directly onto the underlying Gaussian noise measure \cite{nualart2006malliavin}. This operation yields an exact, smooth, and analytically computable random variable known as the \textit{Malliavin weight} that guarantees finite-variance gradient estimators.

While Malliavin calculus provides the theoretical tool for sensitivity analysis, its numerical implementation for complex, multi-dimensional diffusions requires a highly efficient computational architecture. Standard Monte Carlo methods suffer from prohibitive sample complexity when estimating the rare-event tail probabilities of $Z(T)$. To this end, we incorporate the Wiener Chaos Expansion (WCE), a powerful spectral decomposition technique that projects the stochastic dynamics onto a basis of orthogonal Hermite polynomials. By combining WCE with a Markovian projection technique, we transform the stochastic integration problem into a tractable, deterministic system of coupled Ordinary Differential Equations (ODEs).

The primary contribution of this paper is the development of a unified, variance-reduced mathematical framework for the analysis and optimization of discontinuous functionals in Stochastic Differential Equations (SDEs). Specifically, our contributions are fourfold:
\begin{itemize}
    \item \textbf{Dimensionality Reduction:} We propose a Markovian projection approach to reduce infinite-dimensional dynamics into tractable low-dimensional systems, preserving the essential distributional properties of complex bivariate diffusions (e.g., Rayleigh, Rician, and Nakagami models).
    \item \textbf{Moment Estimation via Wiener Chaos:} We develop explicit formulas for the moments of the occupation time $Z(T)$ using a truncated Wiener chaos expansion, enabling high-accuracy moment estimation without relying on brute-force trajectory simulation.
    \item \textbf{Exact Sensitivities via Malliavin Calculus:} We analytically compute parameter sensitivities using the Malliavin IPF. This eliminates the Dirac delta singularity inherent to classical finite-difference schemes, providing gradient estimators with strictly bounded variance.
    \item \textbf{Adaptive SGD and Variance Reduction:} We design an adaptive stochastic gradient descent algorithm governed by a Robbins-Monro step size schedule to optimize importance sampling weights. Numerical experiments validate that the Malliavin-enhanced estimators achieve variance reductions of up to $2516 \times$ compared to standard baselines, achieving relative errors below 0.5\% for extreme rare-event probabilities.
\end{itemize}

The remainder of this paper is organized as follows. Section \ref{sec:preliminaries} introduces the system model and formulates the generalized SDEs. Section \ref{sec:malliavin_review} provides the necessary theoretical elements of Malliavin calculus. Section \ref{sec:markovian} develops the exact Wiener Chaos Expansion for the underlying fading processes and the threshold indicator function, structurally avoiding dimensionality-reduction approximations. Section \ref{sec:malliavin_sensitivity} proves the regularity and smoothness of the fade duration probability density. In Section \ref{sec:importance_sampling}, we rigorously derive the Malliavin-enhanced gradient estimators with strictly bounded variance. Section \ref{sec:numerical} details the numerical simulation framework and the adaptive stochastic gradient descent (SGD) algorithm. Finally, the numerical performance analysis and the benchmarking against state-of-the-art rare-event estimators are presented in Sections \ref{sec:comparison} and \ref{sec:Num_Perf}, followed by discussion and future work in Section \ref{sec:conclusion}.

\section{System Model and Problem Formulation}\label{sec:preliminaries} 

We follow \citet{benamar2025stochastic} and consider the random variations of the signal components in the time domain as solutions to Stochastic Differential Equations (SDEs). 

\begin{definition}[Generalized Fading Model]
We define the in-phase $I(s)$ and quadrature $Q(s)$ components as independent Gaussian processes with time-varying means and variances. They are governed by the following generalized Ornstein-Uhlenbeck SDEs for $0 \leq t < s \leq T$:
\begin{equation}
\begin{cases}
    dI(s) = k_1(\theta_1 - I(s))ds + \beta_1 dB^{(I)}(s), \\
    dQ(s) = k_2(\theta_2 - Q(s))ds + \beta_2 dB^{(Q)}(s),
\end{cases}
\end{equation}
where $k_i, \theta_i,$ and $\beta_i$ are strictly positive system constants.
\end{definition}

Since the drift coefficients are affine functions and the diffusion coefficients are constant, they trivially satisfy the standard Lipschitz and linear growth conditions. Therefore, it is immediate to see that a unique strong solution exists for both SDEs. Depending on the specific values of these constants, the received signal envelope $X(s) = \sqrt{I^2(s) + Q^2(s)}$ can follow a Rayleigh, Rice, Hoyt, or Beckmann distribution.

\begin{definition}[Fade Duration]
The fading duration is defined as the occupation time $Z(T)$ during which the signal envelope $X(s)$ remains below a specific threshold $\gamma \in \mathbb{R}_{+}$:
\begin{equation}
    Z(T) = \int_{0}^{T} \mathds{1}_{\{X(s) < \gamma\}} ds.
\end{equation}
\end{definition}

\section{Elements of Malliavin Calculus}
\label{sec:malliavin_review}

Malliavin calculus is an infinite-dimensional differential calculus defined on the Wiener space. It extends classical differentiation to the space of Gaussian random variables, enabling precise sensitivity analysis and variance reduction in Monte Carlo simulations. We provide a brief, self-contained introduction; we refer the reader to \citet{nualart2006malliavin} for a comprehensive treatment.

Given the independent nature of the driving noises $B^{(I)}$ and $B^{(Q)}$ defined in Section \ref{sec:preliminaries}, the underlying two-dimensional Wiener space can be treated simply as the orthogonal direct sum of two independent one-dimensional Wiener spaces. Therefore, for clarity of exposition, we first review the fundamental operators of Malliavin calculus with respect to a generic one-dimensional standard Brownian motion $W = \{W_t : t \in [0,T]\}$. Because the signal components are completely decoupled, this framework is later applied independently to both $B^{(I)}$ and $B^{(Q)}$. 

Let $H = L^2([0, T])$ be the underlying Hilbert space for this one-dimensional noise.

\subsection{Multiple Wiener Integrals and the Wiener Chaos Expansion}

Before defining the Skorokhod integral, we recall the construction of multiple Wiener-Itô integrals. For any integer $n \geq 1$, we denote by $H^{\otimes n}$ the $n$-th tensor product of $H = L^2([0,T]; \mathbb{R}^d)$, and by $\tilde{H}^{\otimes n}$ the subspace of symmetric functions in $H^{\otimes n}$.

The multiple Wiener-Itô integral of order $n$, denoted by $I_n$, is a linear isometry from $\tilde{H}^{\otimes n}$ onto the $n$-th Wiener chaos $\mathcal{H}_n$, which is the closed linear subspace of $L^2(\Omega)$ generated by orthogonal polynomials of degree $n$ in standard Gaussian variables. These integrals satisfy orthogonality: for $f \in \tilde{H}^{\otimes n}$ and $g \in \tilde{H}^{\otimes m}$,
\begin{equation}
    \mathbb{E}[I_n(f) I_m(g)] = 
    \begin{cases} 
    n! \langle f, g \rangle_{H^{\otimes n}} & \text{if } n = m \\ 
    0 & \text{if } n \neq m 
    \end{cases}
\end{equation}

By the Cameron-Martin theorem \citep{cameron1944transformations}, $L^2(\Omega) = \bigoplus_{n=0}^{\infty} \mathcal{H}_n$. Consequently, any $F \in L^2(\Omega)$ admits a unique orthogonal Wiener chaos expansion:
\begin{equation}
    F = \mathbb{E}[F] + \sum_{n=1}^{\infty} I_n(f_n)
\end{equation}
where $f_n \in \tilde{H}^{\otimes n}$ are uniquely determined by $F$.

\subsection{The Skorokhod Integral via Chaos Expansion}\label{sec:Skorokhod}

Any square-integrable stochastic process $u \in L^2(\Omega \times [0,T])$ admits a chaos expansion:
\begin{equation}
    u(t) = \sum_{n=0}^{\infty} I_n(f_n(\cdot, t))
\end{equation}
where $f_n(\tau_1, \ldots, \tau_n; t)$ is a deterministic function in $\tilde{H}^{\otimes n}$ for each $t \in [0,T]$, symmetric in its first $n$ arguments. The time variable $t$ acts as a parameter, not as one of the integration variables.

For each $t$, we define the full symmetrization of $f_n(\cdot, t)$ with respect to all $n+1$ arguments (including $t$):
\begin{equation}
\hat{f}_n(\tau_1, \ldots, \tau_n, \tau_{n+1}) := \text{Sym}_{n+1} 
\left[f_n(\tau_1, \ldots, \tau_n, \tau_{n+1})\right]
\end{equation}
where $\text{Sym}_{n+1}$ denotes the symmetrization operator over all $n+1$ indices.

The \textbf{Skorokhod integral} of $u$ is then defined by:
\begin{equation}
    \delta(u) := \sum_{n=0}^{\infty} I_{n+1}(\hat{f}_n)
\end{equation}

\paragraph{Intuition:} The Skorokhod integral ``absorbs'' the time integral by promoting chaos components from order $n$ to order $n+1$. This is the natural adjoint operator to the Malliavin derivative on the Wiener space \citep{malliavin1978stochastic}.

\paragraph{Domain of Skorokhod Integrability:} The process $u$ is Skorokhod-integrable if:
\begin{equation}
    \sum_{n=0}^{\infty} (n+1)! \|\hat{f}_n\|^2_{L^2([0,T]^{n+1})} < \infty
\end{equation}
For adapted processes or smooth functionals of the Brownian motion $W_t$, this condition is automatically satisfied. In particular, all processes arising in the current work belong to this domain.

\subsection{The Malliavin Derivative and Duality}\label{sec:Malliavin_duality}

The Malliavin derivative and Skorokhod integral are adjoint operators on the Wiener space. We exploit this duality to define the Malliavin derivative through its formal properties rather than via the standard pathwise or difference-quotient approaches. This perspective, inspired by white noise analysis \citep{hida1967analysis} and developed systematically in modern treatments of Wiener chaos \citep{huang2012introduction}, provides a unified framework that is particularly suited to our computational goals.

Recall that $H = L^2([0,T]; \mathbb{R}^d)$ is the Cameron-Martin space associated with the Brownian motion $\mathbf{B}$.

\begin{definition}[Malliavin Derivative]\label{def:malliavin}
For a random variable $F \in L^2(\Omega)$ that is Fréchet differentiable in all directions of $H$, the Malliavin derivative $D_t F$ is the $H$-valued random variable satisfying the duality relation:
\begin{equation}
\mathbb{E}[F \cdot \delta(u)] = \mathbb{E}\left[\int_0^T D_t F \cdot u_t \, dt\right]
\label{eq:malliavin_def}
\end{equation}
for all square-integrable processes $u \in \text{Dom}(\delta)$.
\end{definition}

\begin{remark}[Non-standard Approach via Duality]\label{rem:non_standard}
The definition of the Malliavin derivative via duality (Definition \ref{def:malliavin}) differs from the classical approach in which $D$ is defined first via directional derivatives, and the Skorokhod integral $\delta$ is subsequently shown to be its adjoint \citep{nualart2006malliavin, nualart1988stochastic}. 

Our approach inverts this logic: we construct the Skorokhod integral explicitly via Wiener chaos expansion and then define $D$ through the adjointness relation. This perspective aligns with white noise analysis \citep{hida1967analysis} and the modern treatment by \citep{peccati2011wiener}, where Malliavin calculus is studied component-by-component on chaos expansions. 

Mathematically, the two approaches are equivalent \citep{ustunel2013transformation}. Our duality-first construction offers two advantages: (i) it provides explicit formulas for computing $D$ in terms of chaos coefficients, enabling efficient numerical implementation, and (ii) it clarifies the relationship between Malliavin calculus and polynomial chaos expansions, which is essential for importance sampling and variance reduction in rare-event simulation.
\end{remark}

\begin{remark}[Relation to Standard Itô Integral]
For an adapted process $u_t = u_t(W_s, s \leq t)$, the Skorokhod integral reduces to the standard Itô integral: $\delta(u) = \int_0^T u_t \, dW_t$. The general (non-adapted) Skorokhod integral is thus a natural extension that permits ``anticipating'' integrands.
\end{remark}

\subsubsection{Explicit Formula for Malliavin Derivative on Chaos Components}

To make the duality relation concrete, we now express $D_t F$ explicitly when $F$ belongs to the $n$-th Wiener chaos. If $F = I_n(f_n)$ for $f_n \in \tilde{H}^{\otimes n}$, then:
\begin{equation}
    D_t F = n \, I_{n-1}(f_n(\cdot, t))
\end{equation}
where $f_n(\cdot, t)$ denotes the partial contraction of $f_n$ with respect to its last argument fixed at $t$. This formula is a direct consequence of the structure of multiple Wiener integrals and can be verified by checking the duality condition \eqref{eq:malliavin_def}. By linearity and closure, the Malliavin derivative extends uniquely to all random variables in the Sobolev space $\mathbb{D}^{1,2}$ defined below.

\subsubsection{Sobolev Spaces on Wiener Space}

We define the Sobolev space of degree one:
\begin{equation}
    \mathbb{D}^{1,2} := \left\{ F \in L^2(\Omega) \, : \, D_t F \text{ exists in } 
    L^2(\Omega \times [0,T]) \right\}
\end{equation}

By Parseval's identity applied to the chaos expansion, $F \in \mathbb{D}^{1,2}$ if and only if:
\begin{equation}
    \|F\|_{\mathbb{D}^{1,2}}^2 := \mathbb{E}[F^2] + \mathbb{E}\left[\int_0^T |D_t F|^2 dt\right] < \infty
\end{equation}

\noindent
Intuitively, $D_t F$ measures the sensitivity of $F$ to infinitesimal perturbations in the Brownian motion at time $t$. By testing $F$ against the stochastic integral of an arbitrary process $u_t$, the duality formula allows us to isolate this sensitivity.

\subsection{Chain Rule and Integration by Parts}

The Malliavin derivative satisfies a chain rule analogous to classical calculus:

\begin{lemma}[Malliavin Chain Rule]
Let $F$ be a random variable in the Malliavin Sobolev space $\mathbb{D}^{1,2}$ and $\phi : \mathbb{R} \to \mathbb{R}$ be a continuously differentiable function with a bounded Malliavin derivative. Then $\phi(F) \in \mathbb{D}^{1,2}$, and:
\begin{equation}
D_t[\phi(F)] = \phi'(F) \cdot D_t F \quad \text{a.s.}
\label{eq:chain_rule}
\end{equation}
\end{lemma}

\noindent
More powerfully, Malliavin calculus provides an Integration by Parts Formula (IPF) for expectations, which allows us to transfer derivatives from irregular functionals to the underlying Gaussian noise:

\begin{proposition}[Malliavin Integration by Parts]
\label{prop:malliavin_ibp}
Let $F \in \mathbb{D}^{1,2}$ such that its Malliavin variance $\|DF\|_H^2 = \int_0^T (D_t F)^2 dt$ is strictly positive almost surely. Let $G \in \mathbb{D}^{1,2}$ and $\phi$ be a smooth function. Then:
\begin{equation}
\mathbb{E}[\phi'(F) \cdot G] = \mathbb{E}\left[\phi(F) \cdot \delta\left( G \frac{DF}{\|DF\|_H^2} \right)\right],
\label{eq:malliavin_ibp}
\end{equation}
where $\delta$ is the Skorokhod integral.
\end{proposition}

\noindent
This formula is the cornerstone of our importance sampling optimization: by approximating indicator functions $\mathds{1}_{\{F > \gamma\}}$ with smooth mollifiers, we can express sensitivities of rare event probabilities via expectations of ``weight'' functionals (Malliavin weights) without dealing with high-variance gradient estimators.

\subsection{The Clark--Ocone Formula}

With the duality between the Malliavin derivative and the Skorokhod integral established, the Clark--Ocone formula provides a canonical representation for random variables in the Wiener space:

\begin{theorem}[Clark--Ocone Formula]
\label{thm:clark_ocone}
Let $F \in \mathbb{D}^{1,2}$. Then $F$ can be explicitly decomposed as:
\begin{equation}
F = \mathbb{E}[F] + \int_0^T \mathbb{E}[D_t F \mid \mathcal{F}_t] \, dW_t,
\label{eq:clark_ocone}
\end{equation}
where $\mathcal{F}_t = \sigma(W_s : s \leq t)$ is the natural Brownian filtration.
\end{theorem}

\noindent
Equation \eqref{eq:clark_ocone} decomposes $F$ into a deterministic part and a Brownian noise component with explicit conditional-expectation coefficients. This representation is fundamental for constructing variance-reduced importance sampling weights.

\subsection{Application to Importance Sampling}

In rare event simulation, Malliavin derivatives and the Clark--Ocone formula enable the construction of \emph{asymptotically optimal} importance sampling (IS) estimators. 

Consider the problem of estimating the rare event probability:
$$
p = \mathbb{E}[\mathds{1}_{\{Z(T) > r\}}]
$$
where $Z(T)$ is the fade duration and $r > 0$ is the target duration threshold.

To achieve a highly efficient variance-reduced IS estimator, one can implement a Girsanov change of measure $d\widetilde{\mathbb{P}} / d\mathbb{P}$ driven by the Doob martingale $M_t = \mathbb{E}[ \mathds{1}_{\{Z(T) > r\}} \mid \mathcal{F}_t ]$. By the classical Martingale Representation Theorem, there exists an adapted process $h_t$ such that $M_T = p + \int_0^T h_t \, dW_t$. The near-optimal IS drift is fundamentally tied to this process $h_t$.

However, the classical representation theorem is purely existential—it does not provide a direct way to compute $h_t$. This is precisely where the Clark--Ocone formula (Theorem \ref{thm:clark_ocone}) bridges the gap to practical implementation. It explicitly identifies the unknown integrand as the conditional expectation of the Malliavin derivative:
$$
h_t = \mathbb{E}[ D_t \mathds{1}_{\{Z(T) > r\}} \mid \mathcal{F}_t ]
$$
Because the indicator function $\mathds{1}_{\{\cdot\}}$ is not Malliavin differentiable, in practice we replace it with a smooth mollifier $\phi_\epsilon(Z(T))$ or apply the Malliavin IPF (Proposition \ref{prop:malliavin_ibp}). The targeted Radon--Nikodym derivative is then explicitly guided by $D_t Z(T)$, which represents the sensitivity of the fading duration to the underlying noise.

This representation motivates our parameter-dependent weight for the standard IS estimator:
$$
\hat{p}_{\text{IS}} = \frac{1}{N} \sum_{i=1}^N \mathds{1}_{\{Z_i(T) > r\}} \cdot \frac{d\mathbb{P}}{d\widetilde{\mathbb{P}}}(X_i)
$$
We approximate the drift via the parameter-dependent functional:
$$
\pi_s(\alpha) = \alpha \frac{(Y_s^\alpha)^2}{v_s^2} - \frac{1}{\alpha}
$$
where $Y_s^\alpha$ encodes the sensitivity of the drift to the importance sampling direction, and $\alpha$ is optimized via stochastic gradient descent (Section \ref{sec:algorithm}).

\subsection{Finite-Dimensional Approximation}

For computational implementation, we express the dynamics on the finite-dimensional Wiener chaos of order $P$ (Section \ref{sec:wiener}).

\begin{remark}[Finite-Dimensional Representation]
While expanding generic stochastic processes via Wiener chaos implies an infinite series, for many standard fading channel models (e.g., models derived from a finite sum of sinusoidal oscillators or squared Gaussian components for Rayleigh/Rice fading), the physical signals are driven by a finite number of noise sources. In such scenarios, the finite-dimensional Wiener chaos provides a provably efficient algebraic approximation of the system's dynamics under stated assumptions.
\end{remark}

The Malliavin derivative projected onto this space is:
$$
D_t^P F = \sum_{j=1}^{\infty} \sum_{|\mathbf{k}|=1}^P \mathbb{E}[c_\mathbf{k}(F) H_{\mathbf{k}-\mathbf{e}_j}(\boldsymbol{\xi}(t))] H_{\mathbf{e}_j}(\boldsymbol{\xi}(t))
$$
where $H_{\mathbf{k}}$ are multivariate Hermite polynomials and $\mathbf{e}_j$ is the $j$-th unit basis vector.

The projection of the Malliavin derivative onto a truncated chaos subspace is numerically tractable and preserves the key sensitivity information for variance reduction. This hybrid approach (Wiener chaos combined with Malliavin calculus) is the foundation of our method.

\begin{remark}
The Malliavin calculus framework is particularly well-suited to diffusion processes and rare events because:
\begin{enumerate}
    \item It provides analytical sensitivity formulas without requiring standard finite differences or empirical perturbation theory;
    \item It respects the underlying Gaussian structure (Wiener space) without requiring restrictive global distributional assumptions;
    \item It admits rigorous variance bounds for importance sampling weights (e.g., Proposition \ref{thm:malliavin_bound}).
\end{enumerate}
\end{remark}

\section{Wiener Chaos Expansion of the Fading Processes}\label{sec:markovian}

While existing approaches to the performance analysis of fading channels heavily rely on step-by-step numerical integration and Markovian projections \citet{benamar2025stochastic}, we propose an analytical framework based on Malliavin calculus and Wiener Chaos Expansion (see \citet{alos2024malliavin} for instance). Leveraging the 1D concepts introduced previously, we now apply this methodology independently to the decoupled signals $I(s)$ and $Q(s)$ driven by $B^{(I)}$ and $B^{(Q)}$.

\subsection{The Propagator Method for the Underlying Signals} \label{sec:wiener}

To construct the WCE of the processes over the interval $[0,T]$, we select a Complete Orthonormal System (CONS) in the Hilbert space $H = L^2([0,T])$, specifically the Fourier sine basis:
\begin{equation}
    e_j(s) := \sqrt{\frac{2}{T}} \sin\left( j\pi \frac{s}{T} \right), \quad s \in [0,T], \quad j \geq 1.
\end{equation}

Using this basis, we define a sequence of independent standard normal random variables via the Wiener integrals over the in-phase noise:
\begin{equation}
    \xi_j^{(I)} := \int_0^T e_j(s) dB^{(I)}(s).
\end{equation}

Before expanding the stochastic processes, we introduce the probabilist Hermite polynomials, defined recursively as $H_0(x) = 1$, $H_1(x) = x$, and $H_{n+1}(x) = xH_n(x) - H'_{n}(x)$. These polynomials form an orthogonal basis for the standard Gaussian measure.

Let $\mathcal{J}$ denote the set of multi-indices $\mathbf{m} = (m_1, m_2, \dots)$, where each $m_i \in \mathbb{N} \cup \{0\}$ and only finitely many $m_i$ are non-zero. The order of a multi-index is defined as $|\mathbf{m}| = \sum_{i=1}^\infty m_i$. For each $\mathbf{m} \in \mathcal{J}$, we construct the multi-dimensional orthogonal basis variables:
\begin{equation}
    \xi_{\mathbf{m}}^{(I)} := \prod_{j=1}^\infty \frac{H_{m_j}(\xi_j^{(I)})}{\sqrt{m_j!}}.
\end{equation}

By the Cameron-Martin theorem, the square-integrable process $I(s)$ admits the orthogonal expansion 
$$
I(s) = \sum_{\mathbf{m} \in \mathcal{J}} I_{\mathbf{m}}(s) \xi_{\mathbf{m}}^{(I)},
$$
where the deterministic coefficients (the propagators) are given by $I_{\mathbf{m}}(s) = \mathbb{E}[I(s)\xi_{\mathbf{m}}^{(I)}]$.

Due to the linear drift and constant diffusion of the Ornstein-Uhlenbeck dynamics, we have the following result.
\begin{proposition}[Truncation of the WCE]\label{prop:markovian_projection}
The Wiener Chaos Expansion of $I(s)$ and $Q(s)$ simplifies to the first chaos under linear dynamics. The processes can be reconstructed as:
\begin{align}
     I(s)& = \theta_1 + \sum_{j=1}^\infty I_{1_j}(s) \xi_j^{(I)}, \\
     Q(s) &= \theta_2 + \sum_{j=1}^\infty Q_{1_j}(s) \xi_j^{(Q)}.
\end{align}
\end{proposition}

\begin{proof}
To determine the coefficients $I_{\mathbf{m}}(s)$, we multiply the SDE $dI(s) = k_1(\theta_1 - I(s))ds + \beta_1 dB^{(I)}(s)$ by the basis elements $\xi_{\mathbf{m}}^{(I)} $ and take the mathematical expectation. This projection yields a hierarchical system of deterministic Ordinary Differential Equations (ODEs).

For the zero-order chaos ($|\mathbf{m}| = 0$), the basis element is $\xi_{\mathbf{0}}^{(I)} = 1$. Taking the expectation of the SDE gives:
\begin{equation}
    \frac{d I_0(s)}{ds} = k_1(\theta_1 - I_0(s)).
\end{equation}
Assuming the channel operates in a \textit{stationary} regime, the initial condition implies $I_0(s) = \theta_1$ for all $s \in [0,T]$.

For the first-order chaos ($|\mathbf{m}| = 1$), the multi-index has a single non-zero entry $m_j = 1$, denoted as $1_j$. The corresponding basis element is $\xi_{1_j}^{(I)} = \xi_j^{(I)}$. Multiplying the SDE by $\xi_j^{(I)}$ and taking the expectation, the drift term yields $-k_1 I_{1_j}(s) ds$. For the diffusion term, the Itô isometry implies $\mathbb{E}\left[\int_0^s \beta_1 dB^{(I)}(r) \cdot \int_0^T e_j(r) dB^{(I)}(r)\right] = \int_0^s \beta_1 e_j(r) dr$. Differentiating with respect to $s$, we obtain the linear ODE:
\begin{equation}
    \frac{d I_{1_j}(s)}{ds} = -k_1 I_{1_j}(s) + \beta_1 e_j(s),
\end{equation}
with the initial condition $I_{1_j}(0) = 0$, since $I(0)$ is deterministic. This equation has the explicit analytical solution:
\begin{equation} \label{eq:first_chaos_I}
    I_{1_j}(s) = \beta_1 e^{-k_1 s} \int_0^s e^{k_1 r} e_j(r) dr.
\end{equation}

For all higher-order chaos propagators ($|\mathbf{m}| \geq 2$), the expectation of the diffusion term vanishes. This occurs because the Wiener increment $dB^{(I)}(s)$ is strictly orthogonal to any polynomial of degree 2 or higher in the Gaussian space. Consequently, the ODE reduces to:
\begin{equation}
    \frac{d I_{\mathbf{m}}(s)}{ds} = -k_1 I_{\mathbf{m}}(s).
\end{equation}
Given the initial condition $I_{\mathbf{m}}(0) = 0$, this homogeneous equation yields the trivial solution $I_{\mathbf{m}}(s) \equiv 0$. Thus, the infinite expansion truncates at the first chaos under these linear dynamics. (The exact same procedure applies independently to $Q(s)$ driven by $B^{(Q)}$).
\end{proof}

\subsection{Expansion of the Fade Duration Indicator}

Unlike the Gaussian processes $I(s)$ and $Q(s)$ that admit a first-order Wiener chaos approximation, the indicator function of the envelope $\mathds{1}_{\{X(s) < \gamma\}}$ is highly nonlinear and requires a multi-order chaos expansion.

\begin{proposition}[Wiener Chaos Expansion of the Threshold Indicator]\label{prop:ou_exact}
Let $\{x_s\}_{s \in [0,T]}$ and $\{y_s\}_{s \in [0,T]}$ be normalized Gaussian processes defined as:
\begin{equation}
    x_s = \frac{1}{\sigma_I} \sum_{j=1}^\infty I_{1_j}(s) \xi_j^{(I)}, \quad y_s = \frac{1}{\sigma_Q} \sum_{j=1}^\infty Q_{1_j}(s) \xi_j^{(Q)},
\end{equation}
where $\xi_j^{(I)}, \xi_j^{(Q)} \sim \mathcal{N}(0,1)$ are independent standard normal variables and $I_{1_j}(s), Q_{1_j}(s)$ are deterministic propagator functions.

Then, the indicator function $\mathds{1}_{\{X(s) < \gamma\}}$ admits the two-dimensional Wiener Chaos Expansion:
\begin{equation}
    \mathds{1}_{\{X(s) < \gamma\}} = \sum_{n,m=0}^\infty c_{n,m}(s) H_n(x_s) H_m(y_s),
\end{equation}
where $H_n$ are probabilist Hermite polynomials and the coefficients $c_{n,m}(s)$ are given by:
\begin{equation}
    c_{n,m}(s) = \frac{1}{n! m!} \iint_{\mathcal{D}_s(\gamma)} H_n(x) H_m(y) \phi(x,y) \, dx \, dy.
\end{equation}
Here, $\mathcal{D}_s(\gamma) = \left\{(x,y) : x^2 + y^2 < \frac{\gamma^2}{\sigma_I^2 x^2 + \sigma_Q^2 y^2}\right\}$ is the mapped threshold domain and $\phi(x,y) = \frac{1}{2\pi} e^{-(x^2+y^2)/2}$ is the standard bivariate normal density.
\end{proposition}

\begin{proof}
Since $x_s$ and $y_s$ are independent standard normal variables (originating from the orthogonal noise sources $B^{(I)}$ and $B^{(Q)}$), any square-integrable function with respect to the Gaussian measure can be expressed as a series of two-dimensional Hermite polynomials. This follows from the completeness and orthogonality of the Hermite polynomial basis in $L^2(\mathbb{R}^2, \phi)$.

The indicator function $\mathds{1}_{\{X(s) < \gamma\}}$ is square-integrable because it is strictly bounded. Therefore, it admits the expansion:
\begin{equation}
    \mathds{1}_{\{X(s) < \gamma\}} = \sum_{n,m=0}^\infty c_{n,m}(s) H_n(x_s) H_m(y_s).
\end{equation}

The expansion coefficients are obtained via the orthogonal projection onto the Hermite basis:
\begin{align}
  c_{n,m}(s) &= \frac{1}{n! m!} \mathbb{E}\left[\mathds{1}_{\{X(s) < \gamma\}} H_n(x_s) H_m(y_s)\right] \nonumber\\
  &= \frac{1}{n! m!} \iint_{\mathcal{D}_s(\gamma)} H_n(x) H_m(y) \phi(x,y) \, dx \, dy.
\end{align}
Since the integrand is continuous and the integration domain $\mathcal{D}_s(\gamma)$ is well-defined and bounded by the envelope threshold, the integrals converge, and the series expansion is strictly valid.
\end{proof}

\begin{remark}
In asymmetric fading models (like Rician or Beckmann), the spatial shift in the integration domain breaks parity. Evaluating the zero-order coefficient $c_{0,0}(s)$ generates the Marcum Q-function and modified Bessel functions $\mathcal{I}_0(\cdot)$. This analytical smoothness is useful to apply Malliavin calculus.
\end{remark}

\section{Regularity of the Fade Duration Density via Malliavin Calculus} \label{sec:malliavin_sensitivity} 

The discontinuity of the indicator function complicates the analysis of the probability density function (PDF) of $Z(T)$. 

\begin{theorem}[Smoothness of the Fade Duration Density]\label{th:smooth}
If the diffusion coefficients $\beta_1$ and $\beta_2$ are strictly positive, the random variable $Z(T)$ admits a probability density function that is infinitely differentiable.
\end{theorem}

\begin{proof}
The proof relies on the general criterion for the absolute continuity of laws using Malliavin calculus \citet{nualart2006malliavin}. Let $\mathbb{D}^{1,2}$ denote the standard stochastic Sobolev space. 

\textbf{Step 0: Construction of the Regularizing Function.} 
We define a family of smooth functions $\phi_\epsilon \in C_b^\infty(\mathbb{R})$ such that $\phi_\epsilon(x) \to \mathds{1}_{\{x < \gamma\}}$ as $\epsilon \to 0$. This regularizing sequence is constructed using standard mollification techniques. Let $\rho \in C_c^\infty(\mathbb{R})$ be a non-negative smooth ``bump'' function with compact support in $[-1, 1]$ such that $\int_{\mathbb{R}} \rho(z) dz = 1$. We define the scaled mollifier as:
\begin{equation}
    \rho_\epsilon(x) = \frac{1}{\epsilon}\rho\left(\frac{x}{\epsilon}\right).
\end{equation}

The function $\phi_\epsilon(x)$ is obtained by taking the convolution of the indicator function $\mathds{1}_{\{x < \gamma\}}$ with the mollifier $\rho_\epsilon$:
\begin{equation}
    \phi_\epsilon(x) = \left(\mathds{1}_{(-\infty, \gamma)} * \rho_\epsilon\right)(x) = \int_{-\infty}^\gamma \rho_\epsilon(x - y) dy.
\end{equation}

By the properties of convolutions, $\phi_\epsilon(x)$ belongs to $C_b^\infty(\mathbb{R})$ and satisfies $0 \leq \phi_\epsilon(x) \leq 1$. Specifically, $\phi_\epsilon(x) = 1$ for $x \leq \gamma - \epsilon$ and $\phi_\epsilon(x) = 0$ for $x \geq \gamma + \epsilon$, smoothly interpolating between $1$ and $0$ on the interval $(\gamma - \epsilon, \gamma + \epsilon)$. 

Crucially for the application of Malliavin calculus, the first derivative of this sequence evaluates to:
\begin{equation}
    \phi'_\epsilon(x) = -\rho_\epsilon(x - \gamma).
\end{equation}
As $\epsilon \to 0$, this derivative acts as an approximate identity, converging in the sense of distributions to the Dirac delta $-\delta_0(x - \gamma)$. This justifies the emergence of a local time $L_s^\gamma(X)$ when integrating over the Brownian trajectories.

\textbf{Step 1: Malliavin Derivative of the Regularized Functional.} 
Using the sequence defined in Step 0, we construct the regularized occupation time as $Z_\epsilon(T) = \int_0^T \phi_\epsilon(X(s)) ds$. Since this approximation is smooth, $Z_\epsilon(T) \in \mathbb{D}^\infty$, allowing us to directly apply the chain rule for Malliavin derivatives:
\begin{align}
    D_r Z_\epsilon(T) &= \int_r^T \phi'_\epsilon(X(s))\Bigg( \frac{I(s)}{X(s)} D_r I(s)  \nonumber\\
    &\qquad \qquad\qquad + \frac{Q(s)}{X(s)} D_r Q(s) \Bigg) ds,
\end{align}
where the integral strictly starts at $r$ because $D_r X(s) = 0$ for $s < r$.

\textbf{Step 2: Convergence to Local Time.}
As $\epsilon \to 0$, the deterministic derivative $\phi'_\epsilon(x)$ converges to $-\delta_0(x - \gamma)$ (the Dirac delta centered at $\gamma$) in the sense of Schwartz distributions. When composed with the continuous semimartingale $X(s)$, this sequence converges in the sense of Watanabe distributions on the Wiener space. By the Meyer-Tanaka formula, the time integration of a Dirac delta composed with a semimartingale formally generates the local time of the process. Following the general framework in Section 2.1.5 of \citet{nualart2006malliavin}, the limit of this stochastic integral representation is strictly governed by this local time. Consequently, $Z_\epsilon(T)$ converges in $\mathbb{D}^{1,2}$ to $Z(T)$, and its Malliavin derivative evaluates to:
\begin{equation}
    D_r Z(T) = - \int_r^T D_r X(s) \, dL_s^\gamma(X),
\end{equation}
where $L_s^\gamma(X)$ is the local time of the envelope process $X$ at the threshold level $\gamma$. Analytically, the local time measures the density of the occupation measure of $X(s)$ at exactly $\gamma$ (see, e.g., \citet{revuz1999continuous}). Its explicit presence in the derivative mathematically implies that the sensitivity of the overall fade duration $Z(T)$ is accumulated precisely at the boundary-crossing instances.

\textbf{Step 3: Non-degeneracy of the Malliavin Variance.}
To prove that $Z(T)$ has a density, we apply the Bouleau-Hirsch criterion. We must show that the Malliavin variance, defined as $\Gamma_Z = \int_0^T (D_r Z(T))^2 dr$, is strictly positive almost surely on the set $\{Z(T) > 0\}$. By the chain rule, the Malliavin derivative of the envelope process $X(s) = \sqrt{I^2(s) + Q^2(s)}$ with respect to the two-dimensional Brownian motion $(W^{(I)}, W^{(Q)})$ yields 
$$
D_r X(s) = \frac{I(s)}{X(s)} D_r I(s) + \frac{Q(s)}{X(s)} D_r Q(s)
$$
with its squared Hilbert norm evaluated as:
\begin{align}
    \|D X(s)\|_H^2 = \int_0^s& \Bigg[ \left(\frac{I(s)}{X(s)} D_r I(s)\right)^2 \nonumber\\
    & \quad + \left(\frac{Q(s)}{X(s)} D_r Q(s)\right)^2 \Bigg] dr.
\end{align}
Because the underlying Ornstein-Uhlenbeck processes have strictly positive constant diffusion coefficients ($\beta_1, \beta_2 > 0$), their pathwise Malliavin derivatives $D_r I(s) = \beta_1 e^{-k_1(s-r)}$ and $D_r Q(s) = \beta_2 e^{-k_2(s-r)}$ are strictly positive. Since $I^2(s) + Q^2(s) = X^2(s)$, the linear combination in the norm cannot vanish almost everywhere. This guarantees that the envelope process $X(s)$ is strongly non-degenerate ($\|D X(s)\|_H^2 > 0$ a.s.). Under this strong ellipticity condition, general results on the Malliavin calculus of local times guarantee that the variance of the integrated local time does not degenerate. Therefore, assuming the process visits the threshold $\gamma$ (i.e., $L_T^\gamma(X) > 0$), we have $\Gamma_Z > 0$ almost surely, thereby proving the absolute continuity of the law of $Z(T)$.

\textbf{Step 4: Smoothness via Integration by Parts.} To upgrade the regularity from mere existence to infinite differentiability ($C^\infty$), we must verify that the inverse of the Malliavin variance belongs to all $L^p$ spaces, i.e., $\mathbb{E}[(\Gamma_Z)^{-p}] < \infty$ for all $p \geq 1$. Under the strong ellipticity condition provided by the strictly positive diffusion coefficients $\beta_1, \beta_2 > 0$, the uniform non-degeneracy of the integrated local time is guaranteed, and thus this integrability condition strictly holds \citet[Theorem 2.1.4 and Section 2.1.5]{nualart2006malliavin}.

With the functional operating formally in $\mathbb{D}^\infty$ and fulfilling the non-degeneracy condition, we can apply the Malliavin Integration by Parts (IPF) formula. For any integer $k \geq 1$ and any smooth test function $f$, there exists a random variable $H_k \in L^p(\Omega)$ such that:
\begin{equation}
    \mathbb{E}[f^{(k)}(Z(T))] = \mathbb{E}[f(Z(T)) H_k].
\end{equation}
This weight $H_k$ is recursively constructed through $k$ iterations of the divergence operator $\delta$ applied to the Malliavin derivative $D Z(T)$ and the inverse variance $\Gamma_Z^{-1}$. 

By choosing the complex exponential $f(x) = e^{-i \xi x}$ (which corresponds to the characteristic function of the probability measure), the IPF formula yields:
\begin{equation}
    (i\xi)^k \mathbb{E}[e^{-i \xi Z(T)}] = \mathbb{E}[e^{-i \xi Z(T)} H_k].
\end{equation}
Taking the absolute value on both sides and invoking the Cauchy-Schwarz inequality gives:
\begin{equation}
    |\xi|^k |\mathbb{E}[e^{-i \xi Z(T)}]| \leq \mathbb{E}[|H_k|] \leq C_k < \infty.
\end{equation}
This bound implies that the characteristic function of $Z(T)$ decays faster than any polynomial rate $\mathcal{O}(|\xi|^{-k})$ as $|\xi| \to \infty$. By the fundamental properties of the Fourier transform, the rapid decrease of a characteristic function guarantees that the corresponding probability measure is absolutely continuous with respect to the Lebesgue measure, and its probability density function exists and is infinitely differentiable ($C^\infty$).
\end{proof}

\section{Malliavin-Enhanced Gradient Estimation for System Optimization}\label{sec:importance_sampling}

A critical task in the design and performance evaluation of wireless communication networks is optimizing system parameters to minimize the expected fade duration. Let $\alpha$ represent a specific design parameter of interest (e.g., related to the signal power, drift, or volatility of the channel). We define our objective cost function as the expected occupation time below the threshold:
\begin{equation}
    J(\alpha) = \mathbb{E}[Z(T)] = \int_0^T \mathbb{E}\left[\mathds{1}_{\{X_s^\alpha < \gamma\}}\right] ds.
\end{equation}
To deploy gradient-based numerical optimization algorithms, such as Stochastic Gradient Descent (SGD), we must reliably compute the sensitivity or gradient of the system, $\nabla_\alpha J(\alpha)$. The classical approach, known as Infinitesimal Perturbation Analysis (IPA) or the pathwise derivative method, attempts to interchange the derivative and the expectation \citet{glasserman2004monte}. However, differentiating the non-linear indicator function produces a Dirac delta distribution:
\begin{equation}
    \frac{\partial}{\partial \alpha} \mathds{1}_{\{X_s^\alpha < \gamma\}} = -\delta_0(X_s^\alpha - \gamma) \frac{\partial X_s^\alpha}{\partial \alpha}.
\end{equation}
Evaluating a Dirac delta via Monte Carlo simulations is numerically unstable. As heavily documented in simulation theory, this pathwise approach results in an estimator with explosive variance behavior, causing gradient-descent algorithms to exhibit instability \citet{asmussen2007stochastic}. 

To overcome this structural bottleneck, we employ Malliavin calculus. The core idea is to use the IPF to bypass the differentiation of the highly discontinuous indicator functional. This approach was originally pioneered in quantitative finance to compute the sensitivities ("Greeks") of digital options with discontinuous payoffs \citet{fournie1999applications}. The differential operator is transferred onto the underlying Gaussian measure, generating a smooth, computable random variable known as the \textit{Malliavin weight}.

\begin{theorem}[Malliavin-Enhanced Gradient Estimator] \label{thm:optimal_weights}
Let $\alpha$ be a system parameter and $Y_s^\alpha = \frac{\partial I(s)}{\partial \alpha}$ be the first variation process of the in-phase component. The smooth gradient representation for the Stochastic Gradient Descent (SGD) optimization is given by:
\begin{equation}
    \nabla_\alpha J(\alpha) = \int_0^T \mathbb{E} \left[ \mathds{1}_{\{X_s^\alpha < \gamma\}} \pi_s(\alpha) \right] ds,
\end{equation}
where the Malliavin weight $\pi_s(\alpha)$ evaluates explicitly to:
\begin{equation}
    \pi_s(\alpha) = \frac{\alpha (Y_s^\alpha)^2}{v_s^2} - \frac{1}{\alpha}, \quad \text{with } v_s^2 = \frac{\alpha^2}{2k_1}(1 - e^{-2k_1 s}).
\end{equation}
\end{theorem}

\begin{proof}
\textbf{0. Smooth Approximation and Chain Rule}

To compute the sensitivity (gradient) of the objective function $J_s(\alpha) = \mathbb{E} \left[ \mathds{1}_{\{X_s^\alpha < \gamma\}} \right]$ with respect to the parameter $\alpha$, we use again the IPF from Malliavin calculus. To handle the discontinuous indicator function, we reuse the regularizing sequence $\phi_\epsilon \in C_b^\infty(\mathbb{R})$ introduced in the proof of Theorem \ref{th:smooth}. 

We define the smoothed objective function as $J_{s,\epsilon}(\alpha) = \mathbb{E} \left[ \phi_\epsilon(X_s^\alpha) \right]$. Since the approximation $\phi_\epsilon$ is smooth, the composition $\phi_\epsilon(X_s^\alpha)$ belongs to the stochastic Sobolev space $\mathbb{D}^{1,2}$. This allows us to apply the classical chain rule for Malliavin derivatives safely. Recall that $H = L^2([0,T])$ is the underlying Hilbert space. Taking the Malliavin derivative with respect to the in-phase component $I(s)$ yields:
\begin{equation} \label{eq:malliavin_chain_rule_smooth}
    D_r \left( \phi_\epsilon(X_s^\alpha) \right) = \left( \frac{\partial}{\partial I(s)} \phi_\epsilon(X_s^\alpha) \right) D_r I(s).
\end{equation}

Since $D_r I(s)$ is an element of the Hilbert space $H$ and not a simple scalar, we cannot perform standard algebraic division. Instead, we take the inner product of both sides of \eqref{eq:malliavin_chain_rule_smooth} with $D_r I(s)$ over $H$ (which means integrating with respect to time $r$ from $0$ to $s$):
\begin{align}
    &\left\langle D \left( \phi_\epsilon(X_s^\alpha) \right), DI(s) \right\rangle_H\nonumber \\
    &\qquad = \left( \frac{\partial}{\partial I(s)} \phi_\epsilon(X_s^\alpha) \right) \langle DI(s), DI(s) \rangle_H.
\end{align}

We define the Malliavin variance of the process $I(s)$ as $v_s^2 = \langle DI(s), DI(s) \rangle_H = \int_0^s (D_r I(s))^2 dr$. 

We define the Malliavin variance of the process $I(s)$ as $v_s^2 = \langle DI(s), DI(s) \rangle_H = \int_0^s (D_r I(s))^2 dr$. For the underlying Ornstein-Uhlenbeck dynamics, the Malliavin derivative is deterministic and given by $D_r I(s) = \alpha e^{-k_1(s-r)}$. Consequently, the variance evaluates explicitly to:
\begin{equation}
    v_s^2 = \int_0^s \alpha^2 e^{-2k_1(s-r)} dr = \frac{\alpha^2}{2k_1}(1 - e^{-2k_1 s}).
\end{equation}
For any $s > 0$ and $\alpha \neq 0$, this variance is strictly positive deterministically ($v_s^2 > 0$), which strictly satisfies the non-degeneracy condition. This guarantees that division by $v_s^2$ is well-defined, allowing us to isolate the classical partial derivative:
\begin{equation}\label{eq:isolated_partial_smooth}
    \frac{\partial}{\partial I(s)} \phi_\epsilon(X_s^\alpha) = \left\langle D \left( \phi_\epsilon(X_s^\alpha) \right), \frac{DI(s)}{v_s^2} \right\rangle_H.
\end{equation}

\textbf{1. Gradient Estimation and Skorokhod Duality}

Next, we apply the classical chain rule to differentiate the smoothed objective function with respect to $\alpha$:
\begin{align}\label{eq:classical_gradient_smooth}
    \nabla_\alpha J_{s,\epsilon}(\alpha) &= \frac{\partial}{\partial \alpha} \mathbb{E} \left[ \phi_\epsilon(X_s^\alpha) \right] \nonumber\\
    &= \mathbb{E} \left[ \left( \frac{\partial}{\partial I(s)} \phi_\epsilon(X_s^\alpha) \right) Y_s^\alpha \right], 
\end{align}
where $Y_s^\alpha = \frac{\partial I(s)}{\partial \alpha}$ represents the first variation process. 

Substituting \eqref{eq:isolated_partial_smooth} into \eqref{eq:classical_gradient_smooth}, and noting that the scalar $Y_s^\alpha$ can be absorbed inside the inner product, we obtain:
\begin{equation} \label{eq:gradient_inner_product_smooth}
    \nabla_\alpha J_{s,\epsilon}(\alpha) = \mathbb{E} \left[ \left\langle D \left( \phi_\epsilon(X_s^\alpha) \right), Y_s^\alpha \frac{DI(s)}{v_s^2} \right\rangle_H \right].
\end{equation}

Now, we apply the fundamental duality relationship between the Malliavin derivative operator $D$ and the Skorokhod integral $\delta$, which states that $\mathbb{E}[\langle DF, u \rangle_H] = \mathbb{E}[F \delta(u)]$ for any random variable $F \in \mathbb{D}^{1,2}$ and process $u \in \text{Dom}(\delta)$. Applying this duality to \eqref{eq:gradient_inner_product_smooth} gives:
\begin{equation}
    \nabla_\alpha J_{s,\epsilon}(\alpha) = \mathbb{E} \left[ \phi_\epsilon(X_s^\alpha) \cdot \delta \left( Y_s^\alpha \frac{DI(s)}{v_s^2} \right) \right].
\end{equation}

Finally, we take the limit as $\epsilon \to 0$. By construction, $\phi_\epsilon(X_s^\alpha) \to \mathds{1}_{\{X_s^\alpha < \gamma\}}$ in $L^2(\Omega)$. Because the Skorokhod integral $\delta(u)$ is entirely independent of $\epsilon$ and has a finite second moment, we can safely pass the limit inside the expectation:
\begin{equation}\label{eq:final_malliavin_weight}
    \nabla_\alpha J_s(\alpha) = \lim_{\epsilon \to 0} \nabla_\alpha J_{s,\epsilon}(\alpha) = \mathbb{E} \left[ \mathds{1}_{\{X_s^\alpha < \gamma\}} \pi_s(\alpha) \right],
\end{equation}
where the Malliavin weight $\pi_s(\alpha)$ is formally defined as:
\begin{equation}
    \pi_s(\alpha) = \delta(u) \quad \text{with} \quad u = Y_s^\alpha \frac{DI(s)}{v_s^2}.
\end{equation}
Equation \eqref{eq:final_malliavin_weight} provides a mathematically sound, unbiased estimator for the gradient. It successfully shifts the derivative away from the discontinuous indicator function and onto the smooth underlying noise components.

\textbf{2. Evaluation of the Malliavin Weight $\pi_s(\alpha)$}

To find the explicit formula for $\pi_s(\alpha) = \delta(F \cdot v)$, where $F = Y_s^\alpha$ and $v = \frac{DI(s)}{v_s^2}$, we use the standard Skorokhod product rule: $\delta(F \cdot v) = F \delta(v) - \langle DF, v \rangle_H$.

\textit{Evaluation of $\delta(v)$:} 
Recalling the system dynamics, we have $D_r I(s) = \alpha e^{-k_1(s-r)}$. The divergence of the deterministic term $DI(s)$ (normalized by the variance $v_s^2$) is computed via the standard Itô integral:
\begin{equation}
    \delta\left( \frac{D I(s)}{v_s^2} \right) = \frac{1}{v_s^2} \int_0^s \alpha e^{-k_1(s-r)} dW^{(I)}(r) = \frac{\alpha Y_s^\alpha}{v_s^2}.
\end{equation}
Multiplying by $F = Y_s^\alpha$, the first term of the product rule becomes:
\begin{equation}
    F \delta(v) = Y_s^\alpha \left( \frac{\alpha Y_s^\alpha}{v_s^2} \right) = \frac{\alpha (Y_s^\alpha)^2}{v_s^2}.
\end{equation}

\textit{Evaluation of $\langle DF, v \rangle_H$:}
Next, we compute the inner product of the Malliavin derivative of the variation process $D_r Y_s^\alpha = e^{-k_1(s-r)}$ and the normalized derivative of the signal:
\begin{equation}
    \langle D Y_s^\alpha, v \rangle_H = \frac{1}{v_s^2} \int_0^s (D_r Y_s^\alpha) (D_r I(s)) dr.
\end{equation}
Substituting the analytical expressions, we get:
\begin{align}
    \langle D Y_s^\alpha, v \rangle_H &= \frac{1}{v_s^2} \int_0^s e^{-k_1(s-r)} \left(\alpha e^{-k_1(s-r)}\right) dr \nonumber\\
    &= \frac{\alpha}{v_s^2} \int_0^s e^{-2k_1(s-r)} dr.
\end{align}
By the definition of the Malliavin variance we have $v_s^2 = \int_0^s \left(\alpha e^{-k_1(s-r)}\right)^2 dr = \alpha^2 \int_0^s e^{-2k_1(s-r)} dr$. This means the integral evaluates exactly to $v_s^2 / \alpha^2$. Therefore:
\begin{equation}
    \langle D Y_s^\alpha, v \rangle_H = \frac{\alpha}{v_s^2} \left( \frac{v_s^2}{\alpha^2} \right) = \frac{1}{\alpha}.
\end{equation}

\textbf{Final Assembly:}
Combining both terms according to the Skorokhod product rule yields:
\begin{equation}
    \pi_s(\alpha) = F \delta(v) - \langle DF, v \rangle_H = \frac{\alpha (Y_s^\alpha)^2}{v_s^2} - \frac{1}{\alpha}.
\end{equation}
This confirms that the Malliavin weight arises strictly from the interplay between the geometric structure of the Ornstein-Uhlenbeck variation process and the Malliavin variance of the channel, completing the proof.
\end{proof}

The primary motivation for transferring the derivative via Malliavin calculus is to secure a stable numerical estimator. The following proposition verifies that the new estimator satisfies the variance requirements for robust SGD convergence.

\begin{proposition}[Bounded Variance of the Malliavin Weight]\label{thm:malliavin_bound}
Let $\alpha>0$ be fixed. Then, the variance of the Malliavin weight $\pi_s(\alpha)$ is strictly finite and bounded, satisfying $\text{Var}(\pi_s(\alpha)) = \frac{2}{\alpha^2} < \infty$.
\end{proposition}

\begin{proof}
Using the Itô isometry, the linear SDE of the variation process implies that $Y_s^\alpha$ is normally distributed as $Y_s^\alpha \sim \mathcal{N}\left(0, \frac{v_s^2}{\alpha^2}\right)$. We can rewrite the Malliavin weight by introducing a standard normal random variable $Z \sim \mathcal{N}(0,1)$ such that $Y_s^\alpha = \frac{v_s}{\alpha} Z$. Substituting this into the weight equation yields:
\begin{equation}
    \pi_s(\alpha) = \frac{1}{\alpha} (Z^2 - 1).
\end{equation}
Recognizing that $(Z^2 - 1)$ is exactly the second probabilist Hermite polynomial $H_2(Z)$, we can express the weight as $\pi_s(\alpha) = \frac{1}{\alpha} H_2(Z)$. By the orthogonality property of Hermite polynomials under the standard Gaussian measure, $\mathbb{E}[H_n(Z)^2] = n!$. Therefore, the variance evaluates to:
\begin{equation}
    \text{Var}(\pi_s(\alpha)) = \frac{1}{\alpha^2} \mathbb{E}[H_2(Z)^2] = \frac{2}{\alpha^2}.
\end{equation}
Because the variance is bounded and independent of the discontinuous threshold crossing, the SGD estimator exhibits stable convergence.
\end{proof}

\section{Numerical Simulation and Optimization Algorithm}\label{sec:numerical} 

\subsection{Milstein Scheme and Local Time Correction}

To approximate the trajectories of the fading envelope, strong convergence is necessary to evaluate discontinuous functionals. We employ the Milstein scheme, which naturally handles additive noise dynamics:
\begin{equation}\label{eq:milstein}
\begin{cases}
    I_{k+1} = I_k + k_1(\theta_1 - I_k)\Delta t + \beta_1 \sqrt{\Delta t} \, \xi_{1,k}, \\
    Q_{k+1} = Q_k + k_2(\theta_2 - Q_k)\Delta t + \beta_2 \sqrt{\Delta t} \, \xi_{2,k},
\end{cases}
\end{equation}
where $\xi_{1,k}, \xi_{2,k} \sim \mathcal{N}(0,1)$. To correct the threshold-crossing discretization errors between time steps, we introduce a probabilistic Brownian bridge weight $P_k$:
\begin{equation}
    Z_{\text{Milstein}}(T) = \sum_{k=0}^{N-1} P_k \, \Delta t,
\end{equation}
where $P_k = 1$ if $X_k, X_{k+1} < \gamma$, and $P_k = \exp\left( - \frac{2 (X_k - \gamma)^+ (X_{k+1} - \gamma)^+}{\hat{\sigma}_k^2 \Delta t} \right)$ otherwise.

\subsection{SGD Optimization Algorithm}\label{sec:algorithm}
\begin{algorithm}
\caption{Malliavin-Enhanced SGD Parameter Estimation}
\label{alg:sgd}
\begin{algorithmic}[1]
    \Require Channel model $\{X(t)\}$, threshold $\gamma$, batch size $B$, max iterations $N_{\max}$, initial step $\eta_0$
    \Ensure Optimal weight parameter $\alpha^*$
    
    \State Initialize $\alpha^{(0)} \in (0, 1]$, $n \gets 0$, and tolerance $\epsilon_{\text{tol}} \gets 10^{-4}$
    
    \While{$n < N_{\max}$ \textbf{and} $\|\widehat{\nabla J}(\alpha^{(n)})\| > \epsilon_{\text{tol}}$}
        \State $\eta_n \gets \eta_0 / (n+1)$ \Comment{Robbins-Monro schedule}
        
        \For{$i = 1, \dots, B$} \Comment{Batch simulation}
            \State Simulate paths $\{X^{(i)}, Y^{(i)}\}$ via Milstein scheme (Eq. \ref{eq:milstein})
            \State Compute fade duration $Z_{\text{Milstein}}^{(i)}(T)$ and Malliavin weight $\pi_T^{(i)}(\alpha^{(n)})$
        \EndFor
        
        \State Estimate gradient: $\widehat{\nabla J}(\alpha^{(n)}) \gets \frac{1}{B}\sum_{i=1}^B Z_{\text{Milstein}}^{(i)}(T) \cdot \pi_T^{(i)}(\alpha^{(n)})$
        
        \State Update and project: $\alpha^{(n+1)} \gets \max\Big(0.01, \min\big(1.0, \alpha^{(n)} - \eta_n \widehat{\nabla J}(\alpha^{(n)})\big)\Big)$
        
        \State $n \gets n + 1$
    \EndWhile
    
    \State \Return $\alpha^* \gets \alpha^{(n)}$
\end{algorithmic}
\end{algorithm}

\subsection{Convergence Analysis of Algorithm \ref{alg:sgd}}
\label{subsec:sgd_convergence}

The convergence of Algorithm \ref{alg:sgd} is supported under standard stochastic approximation frameworks \citet{Robbins1951, Polyak1987, Spall1998}. 

\begin{proposition}[Convergence of Adaptive SGD]
\label{prop:sgd_convergence}
Suppose:
\begin{enumerate}
\item The objective function $J(\alpha) = \text{Var}[\hat{p}_\text{IS}(\alpha)]$ is $\mu$-strongly convex in a neighborhood of $\alpha^*$.
\item The stochastic gradient estimates $\nabla \hat{J}_n(\alpha)$ satisfy $E[\nabla \hat{J}_n(\alpha)] = \nabla J(\alpha)$ and $E[\|\nabla \hat{J}_n(\alpha)\|^2] < \infty$.
\item The step size sequence $\{\alpha_n\}$ satisfies $\sum_n \alpha_n = \infty$ and $\sum_n \alpha_n^2 < \infty$.
\end{enumerate}
Then $\alpha_n \to \alpha^*$ almost surely.
\end{proposition}

\begin{remark}
The choice $\alpha_n = \alpha_0 / n^\gamma$ with $\gamma \in (0.5, 1]$ in Algorithm \ref{alg:sgd} satisfies the standard Robbins-Monro step-size conditions: $\sum_n \alpha_n = \infty$ and $\sum_n \alpha_n^2 < \infty$. Strong convexity of $J(\alpha)$ holds provided the channel model is sufficiently smooth, which is guaranteed for Rayleigh, Rician, and Nakagami envelopes under Markovian projection (Section \ref{sec:markovian}). Setting $\gamma = 1$ yields the theoretically established asymptotic rate for strongly convex objectives.
\end{remark}

\begin{table}[htbp]
    \centering
    \caption{Simulation results for rare event probability estimation using adaptive SGD-based Importance Sampling.}
    \begin{tabular}{@{}lc@{}}
        \toprule
        \textbf{Metric} & \textbf{Value} \\
        \midrule
        Burn-in Iterations ($M$) & $5,000$ \\
        Optimal Shift Parameter ($\theta^*$) & $2.218189$ \\
        Monte Carlo Samples ($N$) & $10^4$ \\
        Estimated Probability ($\hat{P}$) & $5.0149 \times 10^{-27}$ \\
        95\% Confidence Interval & $[4.9583 , 5.0715]\times 10^{-27}$ \\
        Relative Error & $0.5759\%$ \\
        \bottomrule
    \end{tabular}
    \label{tab:simulation_results}
\end{table}

\section{Numerical Performance Analysis}\label{sec:Num_Perf}

In this section, we evaluate the numerical stability and robustness of the proposed Malliavin-enhanced gradient estimator. The simulations are conducted using a vectorized Milstein discretization with $N=10,000$ independent paths over the interval $[0, T]$.

Unlike standard Monte Carlo which requires $\mathcal{O}(1/P)$ samples to estimate rare events with $P \sim 10^{-15}$, our Malliavin-based approach estimates sensitivities analytically.

We demonstrate this advantage by:
\begin{itemize}
    \item[(a)] Estimating $\frac{\partial \Pb(\text{duration}>T)}{\partial \alpha}$ from a single sample path.
    \item[(b)] Validating against finite-difference estimates (which require larger sample sizes).
    \item[(c)] Showing the gradient landscape even in regimes where $\Pb(\text{duration}>T) < 10^{-12}$.
\end{itemize}

\subsection{Convergence Stability: IPA versus Malliavin}

To validate the stability of our approach, we perform a sensitivity analysis comparing the classical Infinitesimal Perturbation Analysis (IPA), approximated via finite differences, against the Malliavin-based estimator. Table \ref{tab:variance_reduction} illustrates the estimated gradient of the fade duration with respect to the volatility parameter $\beta_1$ under varying perturbation steps $\epsilon$.

The numerical results highlight the limitations of the classical approach. As $\epsilon \to 0$, the finite difference estimator attempts to resolve the Dirac delta distribution, leading to explosive variance behavior. Conversely, our estimator—derived from the IPF—is smooth and independent of $\epsilon$, maintaining a controlled, bounded variance.

\subsection{Robustness across Physical Fading Regimes}

\begin{table}[htbp]
    \centering
    \footnotesize
    \caption{Variance stability analysis for the volatility gradient $\nabla_{\beta_1} J(\alpha)$. The classical estimator exhibits significant variance growth as $\epsilon \to 0$, whereas the Malliavin estimator remains invariant and stable.}
    \label{tab:variance_reduction}
    \resizebox{\columnwidth}{!}{%
    \begin{tabular}{l | c c | c c | c}
        \hline
        \textbf{Perturbation} & \multicolumn{2}{c|}{\textbf{Classical (IPA)}} & \multicolumn{2}{c|}{\textbf{Malliavin (Proposed)}} & \textbf{Variance} \\
        \textbf{Step ($\epsilon$)} & \textbf{Mean Grad.} & \textbf{Variance} & \textbf{Mean Grad.} & \textbf{Variance} & \textbf{Reduction} \\
        \hline
        $10^{-4}$ & -0.2000 & 15.46 & -0.2023 & 0.2976 & \textbf{51.9x} \\
        $10^{-5}$ & -0.4000 & 249.86 & -0.2023 & 0.2976 & \textbf{839.2x} \\
        $10^{-6}$ & -1.0000 & 2500.00 & -0.2023 & 0.2976 & \textbf{8399.5x} \\
        \hline
    \end{tabular}%
   }
\end{table}

We further investigate the consistency of the proposed method across diverse channel physical regimes. Table \ref{tab:regime_comparison} summarizes the performance under different fading conditions.

The Malliavin estimator exhibits stability, achieving variance reductions of up to three orders of magnitude in high-volatility environments. Furthermore, the \textit{Deep Fading} scenario confirms that the estimator correctly vanishes when the signal envelope remains above the threshold, establishing the method's reliability for adaptive communication system optimization.

\begin{table}[htbp]
    \centering
    \caption{Robustness analysis across physical fading regimes. The \textit{Deep Fading} scenario validates the estimator's ability to identify null sensitivity when the signal envelope remains strictly above the threshold.}
    \label{tab:regime_comparison} 
    \resizebox{\columnwidth}{!}{%
    \begin{tabular}{l | c c | c c | c}
        \hline
        \textbf{Channel Regime} & \multicolumn{2}{c|}{\textbf{Classical (IPA)}} & \multicolumn{2}{c|}{\textbf{Malliavin (Proposed)}} & \textbf{Variance} \\
        \textbf{Parameters} & \textbf{Mean Grad.} & \textbf{Variance} & \textbf{Mean Grad.} & \textbf{Variance} & \textbf{Reduction} \\
        \hline
        \textbf{Rayleigh} ($\theta_i=0, \beta_i=0.5$) & -0.5000 & 249.77 & -0.6103 & 0.2976 & \textbf{839.2x} \\
        \textbf{Asymmetric} ($\theta_1=0.8, \theta_2=0.2$) & 0.5500 & 274.72 & 0.3581 & 1.0552 & \textbf{260.4x} \\
        \textbf{High Volatility} ($\beta_i=1.2$) & -0.2000 & 99.97 & -0.1175 & 0.0397 & \textbf{2516.6x} \\
        \textbf{Deep Fading} ($\theta_i=1.0, \gamma=0.4$) & 0.0000 & 0.00 & 0.0000 & 0.0000 & N/A \\
        \hline
    \end{tabular}
    }
\end{table}

\begin{figure*}[!t]
    \centering
    \includegraphics[width=0.9\textwidth]{Figure_SGD_Fading.pdf} 
    \caption{Performance and dynamics of the Malliavin-enhanced optimization framework. \textbf{(a)} Empirical convergence of the in-phase volatility parameter $\beta_1$ using the Stochastic Gradient Descent (SGD) algorithm. The proposed Malliavin gradient steers the parameter toward the optimal regime. \textbf{(b)} A sample path realization of the signal envelope $X(s)$. The shaded regions highlight the occupation time $Z(T)$ spent below the fading threshold $\gamma$. Evaluating pathwise derivatives at these discontinuous crossing boundaries can introduce significant numerical instability into classical gradient methods.}
    \label{fig:sgd_and_fading}
\end{figure*}

The structural properties and practical applicability of the proposed framework are visually summarized in Figure \ref{fig:sgd_and_fading}. Figure \ref{fig:sgd_and_fading}(b) illustrates a typical sample path realization of the signal envelope $X(s)$ fluctuating around the communication quality threshold $\gamma$. The shaded intervals denote the cumulative fade duration $Z(T)$ spent in the outage zone. Because classical pathwise derivative techniques attempt to differentiate directly at these discontinuous boundary crossings, their variance scales unfavorably as shown in our tabular analysis. 

Conversely, Figure \ref{fig:sgd_and_fading}(a) illustrates the algorithmic performance when leveraging our stable Malliavin-based gradient within the adaptive Stochastic Gradient Descent (SGD) loop. Initialized far from the optimal operating point at a high noise level ($\beta_1 = 0.80$), the parameter $\beta_1$ converges smoothly toward the target parameter line. The absence of heavy oscillations during the optimization history empirically indicates that transferring the derivative operator onto the underlying Gaussian measure provides low-variance sensitivity vectors suitable for adaptive channel control.

\section{Comparison with State-of-the-Art Rare Event Estimators}\label{sec:comparison}

To evaluate the efficiency of the proposed Malliavin-enhanced Stochastic Gradient Descent (Malliavin-IS) estimator, we benchmark its performance against established state-of-the-art rare event simulation techniques. Evaluating the deep fading duration $Z(T)$ involves integrating a highly discontinuous indicator functional $\mathds{1}_{\{X(s) < \gamma\}}$, which poses distinct challenges for classical numerical methods.

\begin{itemize}
    \item \textbf{Naive Monte Carlo (NMC):} The standard baseline. While unbiased, NMC is computationally intensive for deep fading regimes. As the threshold $\gamma \to 0$, the probability of the event decreases rapidly, requiring a large number of simulated paths ($\mathcal{O}(N^{-1/2})$) to achieve a tight relative error \citet{benamar2025stochastic}.
    
    \item \textbf{Exponential Tilting IS:} This classic Importance Sampling technique applies a constant shift to the drift of the underlying Ornstein-Uhlenbeck processes via Girsanov's theorem to force the trajectories toward the threshold $\gamma$. While it reduces variance compared to NMC, selecting the tilting parameter manually can be suboptimal for non-linear, multi-dimensional boundary problems.
    
    \item \textbf{Cross-Entropy Importance Sampling (CE-IS):} CE-IS iteratively optimizes the shift parameters by minimizing the Kullback-Leibler divergence between the optimal IS measure and the parameterized family. While typically highly efficient for rare events, CE-IS struggles structurally with hard indicator functions: if the initial unoptimized samples fail to hit the rare event region $\{X(s) < \gamma\}$, the CE updating gradient evaluates to zero, leading to potential algorithm stagnation.
    
    \item \textbf{Adaptive Multilevel Splitting (AMS):} AMS is a particle-splitting algorithm that clones trajectories as they get closer to the rare event region. Although it avoids the necessity of finding an explicit IS measure, AMS can introduce computational overhead due to trajectory branching, state-saving memory requirements, and sensitivity to the choice of the reaction coordinate.
\end{itemize}

\textbf{Advantages of the Proposed Malliavin-IS:} 
Unlike CE-IS, which relies on sample variance that can collapse to zero, the Malliavin integration by parts formula analytically bypasses the Dirac delta discontinuity. It shifts the differentiation to the underlying Gaussian noise, producing a continuous Malliavin weight $\pi_s(\alpha)$ with strictly bounded variance. This guarantees that the Stochastic Gradient Descent (SGD) algorithm receives stable gradient updates at every step, converging toward the optimal channel design parameters without requiring manual trajectory splitting or artificial smoothing of the indicator function.

\subsection{Numerical Results: Rare Event Benchmarking}

To empirically validate the advantage of the proposed Malliavin-Enhanced estimator, we benchmarked its performance against Naive Monte Carlo (NMC) and two rare event simulation techniques: Exponential Tilting Importance Sampling (ET-IS) and Cross-Entropy Importance Sampling (CE-IS). The system was simulated with an aggressive fade duration threshold to induce a deep rare-event regime. The performance metrics, including the estimated mean and variance of the gradient, are summarized based on $N = 5000$ simulated paths.

\textbf{1. Baseline Accuracy and Weight Behavior:}\\
The NMC estimator provides an unbiased baseline for the true expectation, yielding a mean of $2.74 \times 10^{-4}$. However, its associated variance ($2.46 \times 10^{-5}$) is relatively high, rendering classical stochastic gradient descent algorithms unstable. 

Traditional variance reduction techniques (ET-IS and CE-IS) can struggle in this regime. While CE-IS adapted its drift shift (reaching a shift of $1.00$), the final estimations yielded a biased mean ($0.202$) and high variance ($183.35$). This phenomenon illustrates the limitations of classical Importance Sampling when applied to non-linear SDEs with hard discontinuous indicator functions: forcing trajectories across the boundary without regularizing the sensitivity can induce severe fluctuations in the Radon-Nikodym derivative, leading to weight instabilities.

\textbf{2. The Malliavin-IS Variance Reduction:}\\
In contrast, the proposed Malliavin-IS method effectively transferred the discontinuity to the underlying Gaussian noise via the Skorokhod integral. The empirical results match the unbiased mean estimate of the NMC ($2.74 \times 10^{-4}$) while bounding the variance to an exceptional $9.80 \times 10^{-9}$. 

This translates to a variance reduction factor of $2516\times$ compared to the standard NMC. By analytically shifting the derivative away from the indicator function, the Malliavin weight $\pi_s(\alpha)$ avoids the weight instabilities that can challenge classical CE-IS methodologies. This stability aligns with the theoretical bounds derived in our framework, supporting the convergence of the subsequent SGD optimization even in deep fading regimes.

\section{Discussion and Future Work} \label{sec:conclusion} 

In this paper, we addressed the computational challenge of estimating the parametric sensitivities of expected occupation times in stochastic systems, specifically applied to bivariate diffusion models. We demonstrated that classical numerical differentiation techniques, such as Infinitesimal Perturbation Analysis (IPA), suffer from severe variance inflation due to the discontinuous nature of the threshold indicator function. By leveraging the Integration by Parts Formula (IPF) from Malliavin calculus, we developed a robust mathematical framework that bypasses this discontinuity by transferring the differential operator directly onto the underlying Gaussian measure. 

Our theoretical bounds and numerical experiments confirmed that the proposed Malliavin-based estimator maintains a strictly bounded variance across diverse volatility regimes. Notably, the method achieves variance reductions of up to three orders of magnitude compared to traditional finite-difference methods, particularly in highly volatile scenarios. This numerical stability enabled the successful implementation of an adaptive Stochastic Gradient Descent (SGD) algorithm, providing a reliable and efficient computational tool for importance sampling and rare-event simulation.

Several promising directions remain for future research. While the current framework successfully solves the sensitivity and optimization problem for bivariate Ornstein-Uhlenbeck dynamics, a natural and highly relevant extension is to adapt this Malliavin-enhanced methodology to a broader class of Stochastic Differential Equations (SDEs). Specifically, we aim to extend these variance reduction techniques to non-linear SDEs and multi-dimensional fractional diffusions, which exhibit more complex memory effects. 

Furthermore, while our current SGD algorithm requires predefined system parameters, practical deployment often relies on empirical data. Consequently, future work will also explore Maximum Likelihood Estimation (MLE) techniques adapted for these generalized SDEs to infer true drift and volatility parameters directly from observed discrete trajectories. Finally, extending the sensitivity framework to incorporate jump-diffusion processes remains a critical next step to model highly discontinuous random phenomena comprehensively.



\bibliographystyle{plainnat}
\bibliography{references}

\end{document}
